\section{Provenance Model}
\label{sec:provenance}
%%%%%%%%%%%%%%%%%%%%%%%%%%%%%%%%%%%%

More na\"{\i}ve methods of solving \problem{} would execute refined queries on a DBMS directly in order to evaluate their deviation from the constraint set. However, this involves a large amount of redundant computation to compare similar refinements. By capturing the provenance of the query without any filtering, we can use the annotated tuple to evaluate which tuples are present in the top-$k$ for a given refinement. This way we are able to avoid calling into the DBMS in order to evaluate refinements for their deviation from constraint satisfaction. This method is used in \cite{MLJ22} (which itself is inspired by the use of provenance in \cite{DIMT13}) for similar reasons, but requires extension to find refinements in the ranking setting. We therefore define our provenance model similarly in order to extend it here. \\

%%%%%%%%%%%%%%%%%%%%%%%%%%%%%%%%%%%%

\subsection{Provenance Variables \& Expressions}
We annotate tuples with expressions that are used to decide whether the tuple would be matched by a particular query over the data. These expressions are made up of combinations of {\it provenance variables} which produce expressions. 
\begin{definition}[Provenance variables \cite{MLJ22}]
    Given an associated attribute $A$ from $\attr(D)$ and an associated value $v$ from the domain of the attribute $\domain(A, D)$ in a database $D$, a {\it provenance variable} is a variable of the form $A_v$.
\end{definition}
\begin{example}
    For example, a provenance variable generated by the \running{} might be $Ma_{CS}$, where $Ma$ is the attribute `Major' and its associated value is `CS', which is in the domain of `Major'. 
\end{example}
We combine provenance variables into provenance expressions which may be evaluated later in order to determine whether a tuple is selected by a refinement. In our problem, we concern ourselves only with refining the predicates of attributes that are present in the original query. This makes the provenance variables for attributes not in this set irrelevant for deciding whether a tuple is selected, so we are able to restrict the set of provenance variables to attributes only present in the original query. Therefore, we define the set of provenance variables $\provvars$ with respect to a query $Q$ and a database $D$ as the set $\{ A_v \mid \forall A \in \attr(\conds(Q)), v \in \domain(A, D) \}$.
\begin{example}
    The set of provenance variables for the \running{} given the database in \Cref{ex:running} is $\{Ma_{CS}, Ma_{EE},\allowbreak Ma_{ME}, Hr_{110}, Hr_{100}, Hr_{95}, Hr_{80}, Hr_{75}, Hr_{60}\}$, abbreviating {\tt Major} as $Ma$ and {\tt Hours} as $Hr$.
\end{example}

We define provenance variables in order to use them in combination to annotate tuples in the database. In particular, we annotate tuples with the expression $\prov(t)$, which we define as $\sum_{A \in \attr(\conds(Q))} A_{t.A}$. We use these expressions in order to determine which tuples are selected by a given refinement. Deciding this for a tuple given a query is a process called {\it valuation} as described in \cite{MLJ22}, and we define it in the next subsection.
 \begin{example}
 \label{ex:provexpr}
    Let $t$ be the tuple of the first candidate shown in \Cref{tab:candidates}, then $\prov(t) = Ma_{CS} + Hr_{100}$.
\end{example}

%%%%%%%%%%%%%%%%%%%%%%%%%%%%%%%%%%%%

\subsection{Valuation}
Having defined provenance variables and expressions, we now begin to connect them to how they may be used to reevaluate the query with refined predicates without accessing the DBMS for additional query processing. Provenance variables may take on values that are propagated through provenance expressions in a process called {\it valuation} as in \cite{MLJ22}. Essentially, we use the predicates from a query to decide the values of the provenance variables which may then be used to evaluate provenance expressions. We formalize the concept in the definition below.
\begin{definition}[Valuation \cite{MLJ22}]
    Let $A_v$ be a provenance variable. We overload notation to allow evaluating a query $Q$ over $A_v$, as queries are the main source for assigning valuations to provenance variables. Then, $Q(A_v)$ is $1$ if for every predicate $p$ in $\conds(Q)$ such that $p.A = A$, then either $v\ p.\diamond\ p.C$ if $p$ is numerical, or $v \in p.C$ if $p$ is categorical, and is 0 in any other case. We apply this definition recursively through expressions, i.e. $Q(p \diamond p') = Q(p) \diamond Q(p')$ where $p$ and $p'$ are either provenance variables or expressions, and $\diamond$ is some binary operator.
\end{definition}
%Conceptually, we extend the relations of $D$ to conditional tables (c-tables) \cite{IW84}, which associate tuples in the relation with a local condition. We assign the local condition of each tuple $t$ to ${\sf Prov}(t)$.  \resolve{Union/join?}

\begin{example}
    Let $\prov(t)$ be the provenance expression in \Cref{ex:provexpr}, and $Q$ be the \running{}. Then, as we have defined valuation, $Q(\prov(t))$ is $Q(Ma_{CS}) + Q(Hr_{100})$. The valuation of $Q(Ma_{CS})$ is 1, as $Q$ has one predicate over {\tt Major}, and $CS$ is included in its set of constants $C$. The valuation of $Q(Hr_{100})$ is also 1, as $Q$ has one predicate over {\tt Hours} filtering tuples with ${\tt Hours} \geq 90$, and $100 \geq 90$. Finally, $Q(\prov(t))$ is just the sum of these individual valuations, i.e. 2.  
\end{example}

A tuple is only selected by a query if it matches all predicates, or equivalently that all valuations of the provenance variable in its provenance expression are 1. This implies that a tuple $t$ is selected by a query $Q$ if $Q(\prov(t)) \equiv \sum_{p \in \prov(t)} Q(p) = |\attr(\conds(Q))|$. We use this fact to evaluate refinement queries without requiring further involvement of the DBMS. Our model of provenance differs from \cite{MLJ22} in that we do not form provenance expressions from the constraints. The provenance expressions captured consider each tuple independently, and there is no way to extract information about which tuple would come before another in a ranking from just the provenance expressions. Therefore, we are unable to build inequalities for constraints that only consider tuples according to the prefix of the constraint with provenance expressions alone. The next section describes how we use this provenance model while also treating ranking in order to find the refinement that answers \problem{}.  
